\documentclass[12pt]{article}

\usepackage{amsmath, amssymb, amsthm}
\usepackage{enumerate}
\usepackage{hyperref}
\usepackage{geometry}
\usepackage{newtxtext}
\usepackage{newtxmath}
\geometry{margin=1in}

\title{Introduction to Number Theory - Lecture 1}
\author{}
\date{}

% ----------------------------------------------------------------
\begin{document}
\maketitle

I wanted to give a simple and fun problem for everyone that is still quite interesting
and neat to me. This problem is not original and fairly routine.
But, it is a nice example of how to use algebraic techniques to solve a number theory problem.  

We want to find all integer solutions to:
\[
  x^2 + y^2 = z^2.
\]

% ================================================================
\section*{Questions}

\begin{enumerate}[(a)]

  \item Why can we assume without loss of generality that $x$, $y$, and $z$
        are coprime, i.e.\ the greatest common factor of the three integers
        is $1$?

  \item Compute the residue classes of $x^2$, $y^2$, and $z^2$ modulo $4$.
        (\textit{Hint: consider the squares of each residue class mod~$4$.})

        \begin{enumerate}[(i)]
          \item Compute $0^2 \bmod 4$.
          \item Compute $1^2 \bmod 4$.
          \item Compute $2^2 \bmod 4$.
          \item Compute $3^2 \bmod 4$.
        \end{enumerate}

  \item Notice that $x$, $y$, and $z$ cannot all be even.  Assuming $z$ is
        odd, why can we assume $y$ to be odd and $x$ to be even?

  \item State the identity used to rearrange the equation
        $y^2 = (z - x)(z + x)$.  Why are $(z - x)$ and $(z + x)$ coprime? 
        (Hint: compute the sum and difference of $z+x$ and $z-x$.)

  \item Why are $(z - x)$ and $(z + x)$ squares of integers (up to sign)? 
        (Hint: work out some concrete examples and perform prime factorisations)
        Assuming $(z - x) = r^2$ and $(z + x) = s^2$, proceed to the next
        part.

  \item Solve for expressions for $x$, $y$, and $z$ in terms of $r$ and
        $s$. (Hint: compute the sum and difference of $z+x$ and $z-x$.)

  \item Using your expressions from~(f) and the necessary conditions on $r$
        and $s$, compute a concrete Pythagorean triple.

\end{enumerate}

% ================================================================
\newpage
\section*{Solutions}

Credit: R.\ Borcherds,
\href{https://www.youtube.com/watch?v=JZKDmTIFR7A&list}{\textit{YouTube
lecture on Pythagorean triples}}.

\begin{enumerate}[(a)]

  \item We can divide out the greatest common divisor of $x$, $y$, $z$
        throughout: if $d = \gcd(x,y,z)$, then $(x/d,\, y/d,\, z/d)$ is
        again an integer solution, and the new triple is coprime.

  \item Every integer is congruent to $0$, $1$, $2$, or $3$ modulo~$4$.
        \begin{enumerate}[(i)]
          \item $0^2 = 0 \equiv 0 \pmod{4}$.
          \item $1^2 = 1 \equiv 1 \pmod{4}$.
          \item $2^2 = 4 \equiv 0 \pmod{4}$.
          \item $3^2 = 9 \equiv 1 \pmod{4}$.
        \end{enumerate}
        Hence every perfect square is congruent to $0$ or $1$ modulo~$4$.

  \item Since $x^2 + y^2 = z^2$ holds in a commutative ring, and the
        equation is symmetric in $x$ and $y$, we may freely swap $x$ and
        $y$.  Given that $z$ is odd ($z^2 \equiv 1 \pmod 4$), the
        left-hand side must also be $\equiv 1 \pmod 4$, which forces exactly
        one of $x^2,\, y^2$ to be $\equiv 1$ and the other $\equiv 0
        \pmod 4$.  Thus one of $x,y$ is odd and the other even; we label the
        odd one $y$ and the even one $x$.

  \item The factorisation is
        \[
          y^2 = z^2 - x^2 = (z + x)(z - x).
        \]
        Since $\gcd(x,z)=1$ (they are part of a coprime triple), any common
        factor of $(z+x)$ and $(z-x)$ must divide both their sum $2z$ and
        their difference $2x$, hence it divides $\gcd(2z,2x)=2$.  Because
        $z$ is odd and $x$ is even, $z+x$ and $z-x$ are both odd, so $2$
        does not divide them.  Therefore $\gcd(z+x,\, z-x)=1$.

  \item We have two coprime positive integers $(z-x)$ and $(z+x)$ whose
        product is the perfect square $y^2$.  If two coprime integers have a
        product that is a perfect square, each of them must individually be a
        perfect square (up to sign).  Hence $z - x = r^2$ and $z + x = s^2$
        for some positive integers $r$ and $s$.

      Let \(a = z-x\) and \(b = z+x\). Then \(ab = y^2\) and \(\gcd(a,b)=1\).

      We use the key fact:
      \emph{if a prime divides a perfect square, it appears an even number of times.}

      Take any prime \(p\) dividing \(a\). Since \(\gcd(a,b)=1\), the prime \(p\) does not divide \(b\). Hence all occurrences of \(p\) in \(ab=y^2\) come entirely from \(a\).

      But in \(y^2\), every prime occurs an even number of times. Therefore \(p\) occurs an even number of times in \(a\). Since this holds for every prime dividing \(a\), \(a\) is a perfect square. So \(a=r^2\).

      The same argument applies to \(b\), so \(b=s^2\).

      Thus,
      \[
      z-x=r^2,\qquad z+x=s^2.
      \]

  \item Adding and subtracting the two equations $z - x = r^2$ and
        $z + x = s^2$ gives
        \[
          z = \frac{s^2 + r^2}{2}, \qquad
          x = \frac{s^2 - r^2}{2}, \qquad
          y = rs.
        \]
      Let \( r,s \in \mathbb{Z} \) be odd. Then there exist integers \( a,b \in \mathbb{Z} \) such that
      \[
      r = 2a+1, \qquad s = 2b+1.
      \]
      Squaring gives
      \[
      r^2 = (2a+1)^2 = 4a^2 + 4a + 1 = 2(2a^2 + 2a) + 1,
      \]
      so \( r^2 \equiv 1 \pmod{2} \). Similarly,
      \[
      s^2 = (2b+1)^2 = 4b^2 + 4b + 1 = 2(2b^2 + 2b) + 1,
      \]
      so \( s^2 \equiv 1 \pmod{2} \).

      Hence,
      \[
      r^2 + s^2 \equiv 1 + 1 \equiv 0 \pmod{2},
      \]
      so \( r^2 + s^2 \) is even.

      Also,
      \[
      s^2 - r^2 \equiv 1 - 1 \equiv 0 \pmod{2},
      \]
      so \( s^2 - r^2 \) is even as well.

      Therefore both \( r^2 + s^2 \) and \( s^2 - r^2 \) are divisible by \(2\), and hence
      \[
      \frac{r^2 + s^2}{2}, \quad \frac{s^2 - r^2}{2} \in \mathbb{Z}.
      \]

  \item Take $r = 1$, $s = 3$ (odd, $\gcd(r,s)=1$).  Then
        \[
          x = \frac{9-1}{2} = 4, \qquad
          y = 1 \cdot 3 = 3, \qquad
          z = \frac{9+1}{2} = 5.
        \]
        Indeed, $3^2 + 4^2 = 9 + 16 = 25 = 5^2$. \qed

\end{enumerate}

\end{document}
